\chapter{Obiettivi ausiliari}
\label{cap:ausiliari}

Il Capitolo~\ref{cap:grpo} ha mostrato che il reinforcement learning su questo
compito è schiacciato fra gruppi morti e gruppi a varianza nulla. Una risposta
possibile è cercare segnale altrove: non in una reward su sequenze generate, ma
in obiettivi ausiliari calcolabili in modalità \emph{teacher-forced}, cioè sui
riferimenti noti.

Sono stati implementati due obiettivi. Entrambi sono documentati qui insieme ai
criteri che ne governano la promozione a esperimento vero, perché il punto
metodologico del capitolo è proprio questo: \textbf{qualificare un obiettivo
prima di addestrarlo}.

\section{Obiettivo 1: la massa ammessa}

\subsection{Formulazione}

L'idea è di premiare il modello per la probabilità che assegna \emph{all'insieme}
dei token legali, invece che al singolo token corretto. È esattamente la
formulazione~\eqref{eq:allowed-mass} introdotta nei fondamenti:
\begin{equation}
  \label{eq:am-repeat}
  \mathcal{L}_{\text{AM}}
  \;=\;
  -\log \sum_{v \in \allowed_t} \policy(v \mid x, y_{<t}),
\end{equation}
dove $\allowed_t$ è fornito dal Trie del Capitolo~\ref{cap:decoding} tramite il
metodo \texttt{allowed\_mask\_for\_prefixes} e i prefissi $y_{<t}$ sono quelli
del riferimento.

Questo è apprendimento a etichette parziali~\cite{cour2011partial}, equivalente
alla massimizzazione della verosimiglianza marginale~\cite{guu2017mml}. Va
segnalato che nel caso profondo la formulazione ingenua della~\eqref{eq:am-repeat}
è, contrariamente all'intuizione, un punto di partenza competitivo e non un
ripiego~\cite{seohuh2020naive}: è la ragione per cui non è stata introdotta
alcuna variante più elaborata. La loss totale diventa
\begin{equation}
  \mathcal{L} \;=\; \mathcal{L}_{\text{SFT}} \;+\; \lambda\, \mathcal{L}_{\text{AM}}
\end{equation}
con $\lambda$ configurabile.

L'implementazione è in \texttt{src/training/allowed\_mass\_loss.py}, integrata
nel percorso di addestramento tramite
\texttt{src/training/auxiliary\_sft\_trainer.py}.

\subsection{Nota numerica}

Il termine $\log \sum \exp$ implicito nella~\eqref{eq:am-repeat} è calcolato in
precisione a 32 bit, anche quando il resto del passaggio in avanti gira in
precisione ridotta. La ragione è che la somma è su un insieme di migliaia di
elementi con logits di ampiezza variabile, e in precisione dimezzata la perdita
di cifre significative è sufficiente a produrre gradienti sbagliati. È il tipo
di dettaglio che non si nota fino a quando l'addestramento diverge senza motivo
apparente.

\subsection{Che cosa questo obiettivo può e non può fare}

Qui va ripetuta con forza l'osservazione della sezione sui fondamenti: la
~\eqref{eq:am-repeat} è \textbf{piatta} sull'insieme ammesso. Qualunque
ridistribuzione di massa \emph{interna} ad $\allowed_t$ lascia la loss
invariata.

Ne segue una previsione precisa, che è anche il criterio con cui l'obiettivo va
valutato:
\begin{itemize}
  \item l'obiettivo può migliorare la \textbf{calibrazione}: quanta probabilità
        il modello colloca su token legali, e quindi quanta massa il vincolo di
        decodifica deve rimuovere;
  \item l'obiettivo \textbf{non può}, da solo, migliorare l'exact match, perché
        non fornisce alcuna preferenza fra i token legali.
\end{itemize}

Chiunque si aspettasse un miglioramento dell'exact match da questo obiettivo
avrebbe un modello mentale sbagliato della sua forma funzionale. Dichiararlo in
anticipo evita di interpretare un risultato nullo sull'exact match come un
fallimento dell'implementazione.

\section{Obiettivo 2: la loss strutturata}

\subsection{Formulazione}

Il secondo obiettivo introduce un termine che modella le \emph{transizioni} fra
gloss consecutivi, nello spirito dei campi aleatori condizionali (\emph{CRF}) e
delle loro varianti applicate a compiti di etichettatura strutturata
\cite{goyal2019global,huang2021global,zaratiana2023semicrf}.

Il modello è un CRF \emph{sparso a lunghezza esatta}: si assume che la sequenza
di gloss abbia la stessa lunghezza della sequenza sorgente (che è il caso nel
$27{,}5\%$ del corpus, sezione~\ref{sec:artefatti}) e si assegna un punteggio
alle transizioni fra stati consecutivi.

Lo spazio degli stati è ridotto a $513$ elementi: i $512$ gloss più frequenti
più uno stato aggregato \texttt{OTHER} che raccoglie tutti gli altri. Le
transizioni sono stimate esclusivamente dal training set \emph{dopo} la
suddivisione, per evitare contaminazione.

I file coinvolti sono:
\begin{center}
\begin{tabular}{ll}
\toprule
File & Ruolo \\
\midrule
\texttt{src/datasets/structured\_transitions.py} & stima delle transizioni dal training \\
\texttt{src/models/structured\_gloss\_head.py} & la testa strutturata \\
\texttt{src/training/auxiliary\_sft\_trainer.py} & integrazione nella loss \\
\bottomrule
\end{tabular}
\end{center}

\subsection{Vincolo di memoria}

Un CRF ingenuo richiederebbe un tensore di punteggi di forma
$[B, T, |\vocab|]$, con $B$ dimensione del batch, $T$ lunghezza e $|\vocab|$
dimensione del vocabolario. Su un vocabolario di decine di migliaia di elementi
questo è proibitivo.

L'implementazione evita completamente quel tensore usando accumulazione sparsa
(\texttt{scatter\_add}) in precisione a 32 bit: si aggregano i contributi
soltanto per gli indici effettivamente coinvolti. Il requisito di memoria
diventa proporzionale al numero di transizioni osservate, non al prodotto delle
dimensioni.

\subsection{Qualificazione: quanto informa realmente il gloss precedente?}
\label{sec:qualificazione-strutturale}

Prima di spendere tempo di GPU su questo obiettivo, è stata posta la domanda
diretta: \emph{sapere il gloss precedente aggiunge informazione, dato che si
conosce già il token sorgente?}

La domanda è formulabile in modo esatto con l'entropia
condizionale~\eqref{eq:cond-entropy}. La stima è stata condotta su $21\,917$
token esclusi dall'addestramento, con una strategia di \emph{backoff} a tre
livelli
\begin{equation}
  (\texttt{prev}, \texttt{src}) \;\to\; \texttt{src} \;\to\; \text{unigramma}
\end{equation}
per gestire le combinazioni non osservate.

\begin{center}
\begin{tabular}{lr}
\toprule
Quantità & Valore \\
\midrule
$H(\text{gloss} \mid \text{src})$ & $0{,}8560$ bit \\
$H(\text{gloss} \mid \text{src}, \text{prev})$ & $0{,}8116$ bit \\
\midrule
Guadagno & $+0{,}0444$ bit \\
Frazione dell'incertezza residua rimossa & $5{,}19\%$ \\
\bottomrule
\end{tabular}
\end{center}

Il gloss precedente aggiunge $0{,}0444$ bit, cioè rimuove circa il $5\%$
dell'incertezza che resta dopo aver visto il token sorgente. Non è zero: esiste
un prior strutturale reale. Ma è piccolo, e fissa un limite superiore realistico
a quanto un termine di transizione possa contribuire.

\subsection{Il caveat, dichiarato}

La stima soffre di una limitazione che va esposta e non nascosta: è calcolata
solo sulle coppie \emph{allineate}, cioè sul $27{,}5\%$ del corpus in cui testo e
gloss hanno la stessa lunghezza.

Questo è problematico in modo specifico. Il fenomeno linguistico più rilevante
identificato nell'analisi dell'SFT (sezione~\ref{sec:sft-vs-regola-analisi}) è la
cancellazione contestuale della copula \texttt{BE}, responsabile del $65\%$ dei
casi in cui l'SFT batte la baseline a regole. Ma una cancellazione \emph{cambia
la lunghezza}: quei casi stanno per costruzione fuori dal sottoinsieme allineato.

Il valore $+0{,}0444$ bit è dunque una stima ottenuta dove il fenomeno di
interesse non c'è. Potrebbe sottostimare o sovrastimare il vero contributo. È una
limitazione del metodo di stima, non un difetto dell'implementazione, e la sua
risoluzione richiederebbe un allineamento non monotono fra le due sequenze.

\subsection{Onestà sulle ritrattazioni}

Questa stima è stata corretta due volte prima di arrivare al valore riportato.
Le versioni precedenti erano affette da contaminazione fra training e valutazione
e da un backoff mal specificato che assorbiva parte del guadagno apparente. Il
valore $+0{,}0444$ bit è quello sopravvissuto alla verifica.

Segnalarlo non è autodenigrazione: è la condizione perché il numero sia
credibile. Un guadagno di $0{,}0444$ bit è abbastanza piccolo da poter essere
prodotto interamente da un errore di metodo, e chi legge deve sapere che
l'errore è stato cercato.

\section{Il contratto di non regressione}
\label{sec:contratto-ausiliari}

Aggiungere termini alla loss di addestramento è una modifica invasiva: rischia
di alterare il comportamento anche quando è disattivata. Il progetto adotta
quindi un \emph{contratto} verificato da test, il cui contenuto è il seguente.

\begin{description}
  \item[Inerzia esatta a peso nullo.] Con peso $0$ il percorso di addestramento
        deve essere \textbf{bit-identico} a quello del \texttt{SFTTrainer}
        originale. Non ``approssimativamente uguale'': identico. La verifica è
        automatica, e serve a garantire che tutte le celle sperimentali che non
        usano gli obiettivi ausiliari non ne siano influenzate.
  \item[Validazione prima del modello.] La funzione
        \texttt{resolve\_auxiliary\_config} valida la configurazione
        \emph{prima} che il modello venga caricato. Un errore di configurazione
        deve costare secondi, non minuti (si veda la
        sezione~\ref{sec:bug-evalonly} per il caso reale in cui questa regola era
        violata).
  \item[Incompatibilità rifiutate esplicitamente.] Le opzioni di
        \emph{packing} (concatenazione di più esempi in una sequenza) e
        \emph{padding-free} sono rifiutate, perché entrambe distruggono la
        corrispondenza fra posizione nel tensore e posizione nella sequenza
        originale, da cui gli obiettivi ausiliari dipendono. Rifiutare è
        preferibile a produrre silenziosamente risultati sbagliati.
  \item[Un solo processo.] La funzione \texttt{require\_single\_process}
        rifiuta l'esecuzione con più di un processo. La ragione è che
        l'aggregazione delle statistiche strutturate non è stata verificata in
        regime distribuito; piuttosto che dedurre un comportamento non testato,
        l'esecuzione viene bloccata.
  \item[Marcatura delle righe ineleggibili.] Il collatore
        \texttt{CompletionSpanCollator} marca come non eleggibili le righe
        troncate o in cui lo span della completion non è contiguo. Su quelle
        righe gli obiettivi ausiliari non vengono applicati, invece di essere
        applicati a posizioni sbagliate. Una riga è considerata completa se lo
        span supervisionato \emph{contiene} il token di fine sequenza, non se
        \emph{termina} con esso: il template di chat di Qwen chiude il turno
        dell'assistente con \texttt{<|im\_end|>} seguito da un a capo, e il
        controllo sull'ultima posizione rendeva ineleggibile ogni riga.
  \item[Percorsi gold non rappresentabili.] Il grafo delle transizioni è
        stimato solo sul training, con un taglio sparso: un esempio di
        valutazione può usare una transizione che il grafo non contiene. Per
        quella riga il punteggio del percorso gold vale $-\infty$ e la NLL
        $+\infty$. Queste righe sono escluse dalla media e contate a parte
        (\texttt{structured\_unsupported\_rows}), con la stessa politica delle
        righe non mappabili: nessuna riga viene valutata contro un bersaglio
        che il modello non può rappresentare. Prima della correzione una sola
        riga del genere rendeva infinita la loss dell'intero batch.
\end{description}

Il filo comune è: \emph{fallire in modo rumoroso e presto}, invece di produrre
numeri plausibili e sbagliati. È la lezione ricorrente di questo progetto, e il
Capitolo~\ref{cap:infrastruttura} ne raccoglie sette casi concreti.

\section{I criteri di promozione}

Nella prima stesura di questa relazione entrambi gli obiettivi erano
implementati e verificati ma non addestrati, e i criteri qui sotto erano stati
fissati \emph{prima} di spendere tempo di GPU. Nel frattempo entrambi sono stati
addestrati: i risultati sono nel Capitolo~\ref{cap:risultati},
sezione~\ref{sec:osservazione-ausiliari}. I criteri restano quelli dichiarati in
anticipo, ed è su di essi che i risultati vanno letti.

\begin{description}
  \item[Criterio 1 (massa ammessa).] Una sonda in inferenza deve mostrare una
        riduzione misurabile della \emph{massa rimossa} dal vincolo di
        decodifica. Cioè: il modello addestrato con l'obiettivo deve collocare
        più probabilità su token legali, e quindi il Trie deve dover intervenire
        meno. Questa è la previsione che segue dalla forma funzionale
        dell'obiettivo, e se non si verifica l'obiettivo non sta facendo ciò che
        si crede.
  \item[Criterio 2 (loss strutturata).] Un confronto controllato fra la testa
        strutturata addestrata con le transizioni \emph{vere} e la stessa testa
        addestrata con transizioni \emph{mescolate casualmente}. Se le due
        condizioni danno lo stesso risultato, il guadagno non viene
        dall'informazione strutturale ma dalla capacità aggiuntiva della testa, e
        l'obiettivo va abbandonato.
\end{description}

Il secondo criterio è un controllo ablativo classico, e ha una proprietà utile:
è \emph{poco costoso} rispetto all'addestramento completo, perché richiede solo
la testa e non la messa a punto dell'intero modello.

\subsection{Che cosa va mescolato nel controllo}
\label{sec:controllo-mescolato}

``Transizioni mescolate casualmente'' non è una specifica sufficiente, e la prima
implementazione lo ha dimostrato. Quella versione permutava le
\emph{destinazioni} degli archi: una transizione $a \to b$ diventava
$a \to \pi(b)$. Il risultato è un grafo in cui i percorsi gold reali non
esistono più, perché usano archi che la permutazione ha spostato. Su quei
percorsi la loss strutturata non è definita, e con l'esclusione delle righe non
rappresentabili (sezione~\ref{sec:contratto-ausiliari}) il termine si spegne:
nell'addestramento del controllo il $98\%$ dei passi registrati ($2\,012$ su
$2\,051$) non aveva \emph{alcuna} riga valutabile. Il controllo si è quindi
addestrato come un SFT puro, e il confronto misurava ``termine strutturato
contro nessun termine'', non ``transizioni vere contro transizioni mescolate''.

La versione corretta conserva il \emph{supporto} del grafo, cioè l'insieme
esatto delle coppie $(a, b)$, e permuta solo i pesi, separatamente per ogni stato
di partenza. Ogni stato mantiene gli stessi successori ammessi, la stessa
normalizzazione e lo stesso multinsieme di probabilità uscenti (quindi la stessa
entropia); l'unica informazione distrutta è \emph{quale} successore è probabile.
I percorsi gold restano valutabili per costruzione, e il termine strutturato è
attivo con la stessa frequenza del grafo vero.

Il caveat, inevitabile per qualunque controllo su cui la loss sia definita: il
supporto stesso viene dai dati reali. Il controllo isola quindi il valore dei
\emph{pesi} delle transizioni, non quello del supporto.

L'impostazione generale è che un obiettivo ausiliario non va promosso perché è
implementato e i test passano. Va promosso quando esiste una previsione
quantitativa derivata dalla sua forma funzionale, e quella previsione è
verificata su un esperimento di controllo. Su questo progetto, il valore
$+0{,}0444$ bit della sezione~\ref{sec:qualificazione-strutturale} suggerisce che
il margine disponibile sia modesto: è un'informazione che vale la pena avere
\emph{prima} di allocare il calcolo, non dopo.
